\documentclass[11pt]{article}
\usepackage[T1]{fontenc}
\usepackage[utf8]{inputenc}
\usepackage{lmodern}
\usepackage[margin=1in]{geometry}
\usepackage{amsmath,amssymb,amsthm,mathtools}
\usepackage{booktabs,array,longtable}
\usepackage{enumitem}
\usepackage{needspace}
\usepackage{hyperref}
\usepackage{xurl}

\IfFileExists{paper/release_info.tex}{
  \newcommand{\OPHPaperReleaseID}{r2056}
\newcommand{\OPHPaperReleaseDate}{October 31, 2026}
\newcommand{\OPHPaperReleaseBanner}{%
\begin{center}
\small\textbf{Paper release:} \texttt{\OPHPaperReleaseID}\quad\textbf{Released:} \OPHPaperReleaseDate
\end{center}
}
\newcommand{\PaperReleaseID}{\OPHPaperReleaseID}
\newcommand{\PaperReleaseDate}{\OPHPaperReleaseDate}
\newcommand{\PaperReleaseBanner}{%
\begin{center}
\small\textbf{Paper release:} \texttt{\PaperReleaseID}\quad\textbf{Released:} \PaperReleaseDate
\end{center}
}

}{
  \newcommand{\OPHPaperReleaseID}{r2056}
\newcommand{\OPHPaperReleaseDate}{October 31, 2026}
\newcommand{\OPHPaperReleaseBanner}{%
\begin{center}
\small\textbf{Paper release:} \texttt{\OPHPaperReleaseID}\quad\textbf{Released:} \OPHPaperReleaseDate
\end{center}
}
\newcommand{\PaperReleaseID}{\OPHPaperReleaseID}
\newcommand{\PaperReleaseDate}{\OPHPaperReleaseDate}
\newcommand{\PaperReleaseBanner}{%
\begin{center}
\small\textbf{Paper release:} \texttt{\PaperReleaseID}\quad\textbf{Released:} \PaperReleaseDate
\end{center}
}

}

\hypersetup{colorlinks=true,linkcolor=blue,citecolor=blue,urlcolor=blue}
\hypersetup{pdftitle={Standard Photon-Baryon Transport as the Transfer Interface for OPH Sources},pdfauthor={Bernhard Mueller}}
\setlength{\emergencystretch}{3em}
\newcolumntype{L}[1]{>{\raggedright\arraybackslash}p{#1}}

\newtheorem{theorem}{Theorem}[section]
\newtheorem{proposition}[theorem]{Proposition}
\newtheorem{lemma}[theorem]{Lemma}
\newtheorem{corollary}[theorem]{Corollary}
\newtheorem{definition}[theorem]{Definition}
\newtheorem{assumption}[theorem]{Assumption}
\theoremstyle{remark}
\newtheorem{remark}[theorem]{Remark}

\newcommand{\Dd}{\mathrm{d}}
\newcommand{\ah}{\frac{\dot a}{a}}
\newcommand{\kdot}{\dot\kappa}

\title{Standard Photon--Baryon Transport as the\\
Transfer Interface for OPH Sources}
\author{B.\ M\"uller}
\date{\OPHPaperReleaseDate}

\begin{document}
\maketitle
\OPHPaperReleaseBanner

\begin{center}
\small\textbf{Author affiliation:} Bernhard Mueller, Pragma Research Inc.
\end{center}

\begin{abstract}
This note shows that standard general-relativistic photon--baryon Boltzmann
transport is reproduced on an Observer-Patch Holography (OPH) branch that
takes as premises a geometry, photon refinement convergence, an
electromagnetic vertex, a polarized collision matrix, the electron mass, and
a recombination history; none of these is derived in this note. The
reduction gives the linear multipole hierarchy, acoustic and diffusion
approximations, and a line-of-sight transfer formula in a common brightness
convention; these are the Ma--Bertschinger system and its standard
approximations, and nothing in them is derived from observer-patch
consistency. The OPH content is two interfaces, a repair-exchange rate and a
collar stress, which vanish on the standard branch and mark where a
conditional extension would enter.
\end{abstract}

% Shared statement of the OPH tower and its executable hypothesis.
% Included by every active paper. No theorem premise or physical status is
% promoted by this panel; each paper states its own layer and hypothesis.
\paragraph{One theory and its executable hypothesis.}
OPH proposes closure: one complete history is both a timeless
self-referential structure and its own internal computation, and observers
inside it recover and realize its specification. The hypothesis that makes
this proposal concrete and testable is executable.

\begin{quote}\itshape
A simple, fully specified local repair and memory law on the proposed
carrier architecture admits continuing observers whose compatible records
and interventions have the observed \(3+1\)-dimensional quantum-field
description and a common classical-gravity limit. The preferred matter
mechanism is localized topological defects of native overlap transport,
with rest energy derived from their excess maintenance dynamics in the same
physical clock and energy dictionary.
\end{quote}

One source process must meet every part of this statement: records, clocks,
matter, energy and responsive geometry come from the same law, state and
refinement. Coherence, protected codes, radial directions, fields, clocks
and couplings may be specified as part of its design; their implementation
and observable consequences must be established on that process. A failed
specified law rejects that law and its tested interpretation; it rejects the
wider proposal only when the rejected choice is proved necessary.

\paragraph{The tower.}
The papers form one tower over the executable model. Each layer states a
hypothesis and proves finite theorems under displayed premises. A conclusion
passes to the next layer only after that premise is proved on the same
model process; until then it is a conditional implication.
\begin{description}\setlength{\itemsep}{2pt}
\item[0. Closure and agreement.] One history reads itself into existence, and
  agreement on shared records defines the public world. This is the
  interpretation and acceptance criterion; no theorem selects a unique
  universe.
\item[1. Repair mathematics.] Local agreement repair reaches global
  consistency, and nonzero cycle holonomy is the residue that local repair
  cannot remove. Fixed-point, normal-form, frustration and finite
  event-algebra theorems hold on their stated domains.
\item[2. Carrier and repair law.] Federated twelve-port icosahedral carriers
  with the native C/E seam repair are the substrate. The seam rule follows
  from four stated requirements, the quantum record laws are classified, and
  the uniform native repair selects a rank-three slow-response band.
\item[3. Executable observer.] Observers are continuing native processes that
  prepare, read, retain and act on records. Finite complete-output
  executions, overlapping-observer correspondence, event clocks and
  checkpoint continuation are verified; one continuing observer with every
  service and resource accounted is the open join.
\item[4. Clocks and fields.] Modular flow of observer algebras supplies
  candidate local clocks, and native repair supplies a positive continuum
  generator of spectral dimension two on the screen. A declared radial
  extension gives a conditional smooth \(1+3\) scalar description. Complete
  finite source factors join spatial CAR dynamics, particle creation,
  repair and a common calibrated clock. Physical validation of that clock
  and energy normalization, and an interacting chiral \(3+1\) limit, remain
  open.
\item[5. Gauge structure and matter.] Gauge structure is the transport of
  overlap data, and stable massive matter is a localized defect whose
  boundary transport carries a holonomy that admitted repairs cannot erase.
  The compact Lie type \(\mathfrak u(1)\oplus\mathfrak{su}(2)\oplus
  \mathfrak{su}(3)\) follows under complete-response premises, and fixed
  nonzero holonomy forces positive residual mismatch. A specified coherent
  ordering memory protects native exchange-history sectors, and a colored
  extension implements their fermionic exchange. A localized physical
  defect, its maintenance-to-mass law and a predicted measured spectrum
  remain unestablished; the earlier double-load candidate is transient.
\item[6. Gravity.] Einstein dynamics is the coarse response of the same
  process. The null-balance and tensor-completion implication is proved
  under its premises; realizing those premises on one refinement tower is
  open, and the static radial product fails null balance.
\item[7. Cosmology and tests.] The same process must fix the vacuum, dark
  sector, primordial response and black-hole records, and pass independent
  prospective tests. The cosmology papers and the fine-structure, Koide,
  de Sitter and Yang--Mills supplements carry conditional branches,
  diagnostics and retrospective comparisons at this layer.
\end{description}

\paragraph{Completion and stopping.}
The aim is a consistent and complete theory: every layer realized on one
process, or a proved obstruction that removes OPH or a stated branch of it.
Results that no join needs leave the active chain with their provenance. A
universal obstruction over a frozen admissible class, an unavoidable
excluded prediction, or a contradiction among indispensable premises stops
that class. A failed candidate, an unconstructed map or a resource stop has a
narrower meaning. The ordered joins, retained foundations and stopping
rules are recorded with the executable model in the public
repository.\footnote{\url{https://github.com/FloatingPragma/observer-patch-holography/blob/r2056/docs/COMPLETION_PROGRAM.md}}

\paragraph{Fixed physical boundaries.}
Repair counts, host runtime and entropy production are distinct from
energy; the maintenance-to-energy map must be derived with one common
signal speed and no per-species conversion. The modular parameter is
dimensionless until a specified and implemented clock calibration relates
it to duration. Stationary repair or
record entropy is distinct from the fixed-volume generalized-entropy
variation that the Einstein construction consumes. A massless photon makes
the infinite-volume theory gapless, so sector gaps and thresholds name their
state and limit; confinement precludes free massless gluons as observable
particles; a quantum graviton lies outside the classical Einstein theorem.


\section*{This paper's hypothesis and place in the tower}

This note sits at layer seven as the transfer link between a primordial
source and the observed acoustic sky. Its hypothesis is that photon--baryon
transport on an OPH branch is the coarse limit of photon-class records
continued between commits, with Thomson collisions at commits, so that the
branch reduces to the standard Einstein--Boltzmann system plus two OPH
interfaces, a repair-exchange rate and a collar stress, which vanish on the
standard branch.

The reduction is a conditional theorem. The geometry, the photon refinement
convergence condition, the electromagnetic vertex and photon pole, the
polarized collision matrix, the electron mass, the recombination history,
the FLRW background and the initial spectrum are premises; the acoustic,
damping and line-of-sight results are the textbook consequences of the
reduced system. The note certifies no initial spectrum, normalization,
recombination history or likelihood, and it reads no Planck data.

For the reduction to apply, the executable model must realize on one
process a two-polarization photon sector, its continuation between record
commits, a physical clock and a smooth Lorentzian background. The native
construction supplies none of these~\cite{ophmodel,ophspine}, and its
passive repair has stationary decaying correlations and no wave propagator
(Section~\ref{s:premises}). The reduction is refuted on a realized branch
whose photon-class record process has a nonzero scaled quadratic-variation
limit, since the transport is then diffusive or jump-like rather than
Liouville flow, or whose source-level commit kernel differs from the
Thomson functional.

\section{Branch inputs}

\begin{assumption}[Einstein branch]\label{a:geom}
The observer-facing normal form carries a Lorentzian metric $g_{ab}$
satisfying $G_{ab}+\Lambda g_{ab}=8\pi G_{\rm geom}\langle T_{ab}\rangle$,
with $G_{\rm geom}=\ell_\star^2$ from the edge-entropy/area-law theorem and
$\Lambda\ell_\star^2=3\pi/N_\star$ from capacity closure. This area-law
normalization is a premise: a reading of the whole screen's count as the
area of a three-dimensional region's boundary requires a region map that the
executable construction does not supply~\cite{ophentropy}. $N_\star$ is the
capacity-identity estimate read from measured $\Lambda$ through a declared
capacity-to-scale bridge; $\Lambda$ is an input to this assumption, not an
OPH output. A physical interpretation of $N_\star$ beyond that estimate
requires the readback map $F$ and its contraction certificate.
The executable carrier bounds dependence by one seam per ordered repair
event, with no laboratory speed or continuum light cone~\cite{ophcausality},
and does not supply the smooth metric assumed here. This branch requires physically identified
events with exact order reflection or controlled vanishing causal error,
count-to-volume convergence, a common clock and physical stress, together
with the capacity premises. Its smooth Einstein reading additionally needs
same-family tensor-curvature convergence or the continuum
small-ball/null-balance identification; scalar curvature alone is diagnostic.
\end{assumption}

The flat causal-order, count-volume and count-clock limit of the spacetime
construction~\cite{ophEinstein} uses a supplied event law and a relative
reference unit; it does not calibrate the common physical clock required by
the Einstein, electromagnetic and cosmological transfer assumptions.

\begin{assumption}[Electromagnetic action and photon pole]\label{a:qed}
The realized branch carries the Standard Model quotient and the Maxwell
kinetic action with a $\mathrm U(1)$ vertex of strength $\alpha$. This note
uses $\alpha=\alpha_{\rm root}$, the exact root of a declared source map,
conditional on same-scheme transport, target-independent selection of that
map, and a typed bridge showing that its two coordinates read the same
physical quantity. In addition, this transport note
assumes a positive-energy physical QED quantization with two transverse
states and a positive-residue isolated pole at $k^2=0$ on the ordinary
Lorentz-vacuum, unbroken electromagnetic phase. This extra photon-pole
receipt is not implied by the reconstructed abstract group. Charged-lepton
\emph{ratios} are branch outputs; the absolute electron mass $m_e$ is
empirically anchored. The inputs are the cited electroweak branch theorem,
the explicit quantum-pole assumption just stated, and the declared import
$m_e$.
\end{assumption}

\begin{assumption}[Cosmological background]\label{a:bg}
The homogeneous limit of the branch is spatially flat FRW,
$\Dd s^2=a^2(\tau)\,[-\Dd\tau^2+\delta_{ij}\Dd x^i\Dd x^j]$, with Friedmann
dynamics from Assumption~\ref{a:geom}. Scalar perturbations are written in
the conformal Newtonian gauge,
\[
\Dd s^2=a^2(\tau)\left[-(1+2\psi)\,\Dd\tau^2+(1-2\phi)\,\delta_{ij}\Dd
x^i\Dd x^j\right],
\]
with conventions of Ma--Bertschinger (1995). Overdots are $\partial_\tau$.
\end{assumption}

\begin{assumption}[Recombination import]\label{a:rec}
The free-electron history $n_e(\tau)$ (equivalently $x_e(\tau)$) is computed
from standard atomic physics (Saha/Peebles-class) using imported atomic data.
It is in the same import class as the atomic clock unit of the gravity
construction: it contains no CMB-target information, so it is a permissible declared input
under the anti-fitting rule, but it is an import, not a branch output.
\end{assumption}

\begin{assumption}[Conditional initial spectrum]\label{a:init}
Initial curvature perturbations $\zeta(\mathbf k)$ are Gaussian with spectrum
$\Delta^2_\zeta(k)=A_s\,(k/k_\ast)^{\,n_s-1}$. The Gaussian and
statistically isotropic character is the maximum-entropy screen-covariance
theorem. A finite tilt certificate requires a source generator, the weighted
orientation-half identity, and a physical clock receipt. The radial theorem
proves that one-shell inversion is nonunique and supplies physical
source-dilation and cross-covariance tomography as conditional uniqueness
routes. Finite source instantiation is not supplied here. This note consumes
\(A_s\) and \(n_s\) as gated inputs and derives no initial spectrum. The
registered branch \(n_s=1-P/48\) fixes \(P\) at its inherited
\(\alpha\)-calibrated value, and the executable observer construction does
not derive that branch from its joint records~\cite{ophpublicdata}.
\end{assumption}

\section{Conditional Refinement Liouville Transport}

\begin{definition}[One-photon record density]
Conditional on the physical pole in Assumption~\ref{a:qed}, let
$\mathcal E_\gamma$ be its positive-energy photon sector. Between record
commits, a single excitation is a record continued
across patch overlaps. Coarse-graining patch position and overlap momentum
labels defines the phase-space record density $f(x^i,q,\hat n,\tau)$ on the
forward mass shell, where $q$ is comoving momentum magnitude and $\hat n$ the
direction.
\end{definition}

\begin{theorem}[Refinement Liouville]\label{t:liouville}
On a branch with the Maxwell action, quantum photon-pole receipt, and overlap
Lorentz kinematics from Assumptions~\ref{a:qed} and~\ref{a:geom}, let $P_r$ be the quotient-lumpable
free-flow transition kernel at step $h_r$. For every frozen smooth compactly
supported test function $\phi$, set
\[
L_r\phi=\frac{P_r\phi-\phi}{h_r},
\qquad
\Gamma_r(\phi)=
\frac{P_r(\phi^2)-(P_r\phi)^2}{h_r}.
\]
The photon refinement convergence condition consists of compact containment/tightness,
convergence of initial laws, uniform-on-compact generator convergence
$L_r\phi\to X_H\phi$ to the null-geodesic Hamiltonian vector field,
vanishing quadratic variation $\Gamma_r(\phi)\to0$, uniqueness of the
Hamiltonian ODE, refinement and schedule naturality, and vanishing boundary
and phase-space-measure defects. Under this receipt, the
between-commit record flow satisfies, in the refinement limit,
\[
\frac{Df}{\Dd\tau}\;\equiv\;
\frac{\partial f}{\partial\tau}
+\frac{\Dd x^i}{\Dd\tau}\frac{\partial f}{\partial x^i}
+\frac{\Dd q}{\Dd\tau}\frac{\partial f}{\partial q}
+\frac{\Dd\hat n}{\Dd\tau}\cdot\frac{\partial f}{\partial\hat n}
\;=\;C[f],
\]
where the characteristics are the null geodesics of $g_{ab}$ and $C[f]$ is
supported on record commits only. In particular the collisionless flow
preserves the Liouville measure.
\end{theorem}

\begin{proof}[Proof sketch under the convergence condition]
For every test function, the discrete process minus the time integral of
$L_r\phi$ is a martingale whose quadratic variation is controlled by
$\Gamma_r(\phi)$. Tightness gives subsequential limits and vanishing
quadratic variation makes every limiting martingale deterministic. Generator
convergence says that each limit solves the martingale problem for $X_H$,
equivalently the null-geodesic Hamiltonian ODE. Uniqueness identifies the
whole limit, independently of the accepted repair schedule. Thus
$\Dd x/\Dd\tau$, $\Dd q/\Dd\tau$, and $\Dd\hat n/\Dd\tau$ are the
geodesic-flow coefficients. On the background of Assumption~\ref{a:bg},
$\Dd q/\Dd\tau=q[\dot\phi-\ldots]$ has the form used in
\S\ref{s:hierarchy}.

\emph{Measure preservation.} The identified limit is Hamiltonian flow on the
mass shell, so it preserves the canonical Liouville measure. The finite
measure-defect clause checks that the quotient measures converge to that
canonical measure rather than to a different coarse volume form. Accepted repair confluence
only removes schedule dependence of committed endpoints. It supplies neither
an inverse between-commit dynamics nor a phase-space Jacobian. The
orientation doubling of a finite label register is bookkeeping and does not
produce either property.

\emph{Convergence premise.} Photon-sector convergence of the finite patch flow to
the displayed generator and zero-quadratic-variation limits is not supplied here. The abstract Lean
layer proves the corresponding schedule-independence and refinement-naturality
statement in its setting. Section~\ref{s:premises} gives the corresponding
simulator test. Theorem~\ref{t:liouville} is conditional on this convergence
premise.
\end{proof}

\begin{remark}[Scalar continuum on the source family]\label{r:scalar-continuum}
On the conservative golden-record population of the spacetime construction,
a positive source-metric scalar action has a proved finite-time continuum
limit. For $m\ge0$ the finite equation $\ddot q+c^2Lq+m^2q=0$ converges to
$u_{tt}-c^2\Delta u+m^2u=0$ with an explicit detector error bound as
$\epsilon\to0$ and $h/\epsilon^2\to0$, and the canonical kick--drift--kick
time discretization of the same action converges with a full-spectrum bound.
For $m=0$ the limiting characteristics are the null cone $|dx|=c\,dt$. The
read law of the same family has cone radii between $c(1-2h/a)\tau$ and
$c\tau$, so the two limiting cones coincide when the wave speed of the action
is identified with the read-law ratio $a/\Delta$; that identification is a
common-model condition and is not a theorem. The high-frequency limit of the
massless wave equation is the null-geodesic Liouville flow on the flat mass
shell: Wigner measures of its solutions are transported by the Hamiltonian
flow of $|\xi|$ (Lions--Paul 1993). On the flat declared family this supplies
the free-flow generator $X_H$ of Theorem~\ref{t:liouville} for a scalar record
field. It supplies neither the two-polarization photon sector, nor the
quotient lumpability and vanishing quadratic variation of the record
process, nor the characteristics of the perturbed metric of
Assumption~\ref{a:bg}. The finite Maxwell dynamics on the Whitney cone gives
constraint-preserving evolution on a supplied domain and no continuum limit.
The photon refinement convergence of Section~\ref{s:premises} therefore
reduces, on the flat family, to the vector, record-process and
curved-background clauses.
\end{remark}

\begin{corollary}[Collision and diffusion alternatives]
If the finite generator decomposes as $L_r=L_r^{\rm free}+C_r$ and
$C_r\phi\to C\phi$, then the limiting density obeys
\[
\partial_\tau f+\{f,H\}=C^*f.
\]
A nonzero limit of $\Gamma_r$ excludes the deterministic Liouville
conclusion. A diffusion conclusion additionally requires vanishing jumps and
a local second-order limiting generator. A jump limit instead contributes
its jump operator to the transport equation.
\end{corollary}

For example, on two states let $Q\phi(i)=\phi(1-i)-\phi(i)$ and
$P_h=I+hQ$, $0<h\leq1$. For $\phi(i)=i$, the scaled conditional
variance is $\Gamma_h(\phi)=1-h\to1$, while the limiting generator
$Q$ describes rate-one jumps between the two states. Thus nonzero
quadratic variation alone does not identify a diffusion.

\begin{remark}
The transport model places the collision functional at record commits and
uses the independently certified measure-preserving flow between them. Claims
about entropy production therefore refer to the commit terms counted by
$C[f]$. The accepted repair relation and the reversible equilibrium or
between-commit dynamics are separate structures.
\end{remark}

\section{Record Commits and the Thomson Collision Functional}

\begin{theorem}[Scalar-intensity commit-collision functional]\label{t:collision}
Let commits between photon-class and electron-class records occur with the
soft-limit matrix element of the recovered $\mathrm U(1)$ vertex from
Assumption~\ref{a:qed}.
Then, to first order in perturbations and at low occupancy, the collision
functional in Theorem~\ref{t:liouville} is the Thomson functional: in the
electron rest frame,
\[
C[f](q,\hat n)\;=\;a\,n_e\sigma_T\left[
\frac{3}{16\pi}\int \Dd\Omega'\,(1+(\hat n\cdot\hat n')^2)\,f(q,\hat n')
- f(q,\hat n)\right],
\qquad
\sigma_T=\frac{8\pi}{3}\frac{\alpha^2}{m_e^2},
\]
boosted to the frame where baryons move with velocity $v_b$. The
source-derived reversible proposal kernel is assumed to obey detailed balance
on the angular state space. Write/check orientation labels alone do not imply
that property, and the accepted strict-descent repair relation is not the
equilibrium kernel. For every fixed photon energy $q$, isotropic angular
distributions are stationary under this Thomson operator. This displayed
scalar-intensity kernel does not by itself derive the Stokes-space Thomson
matrix needed for polarization.
\end{theorem}

\begin{proof}[Proof sketch]
At low occupancy the commit rate between momentum cells is bilinear in the
occupancies with kernel given by the squared vertex; the soft
(photon energy $\ll m_e$) limit of the recovered $\mathrm U(1)$ vertex is
classical Thomson scattering with the stated total cross-section and
$1+\cos^2$ angular kernel; this step is standard QED evaluated on the branch
couplings and is not repeated here. The detailed-balance receipt equates the
stationary weights of each reversible proposal pair. The resulting symmetric
rest-frame kernel damps angular anisotropy at each fixed
$q$. It does not mix frequencies, so its stationary kernel contains every
isotropic spectrum $f(q)$ rather than a unique Bose or Planck occupancy.
Momentum conservation per commit, summed over
commits, gives equal and opposite momentum transfer densities to the
electron--baryon fluid, with coefficient fixed by
$\bar\rho_\gamma/\bar\rho_b$ bookkeeping; this is the Compton drag used in
\S\ref{s:hierarchy}.
\end{proof}

\begin{remark}[Polarization provenance]
The polarization hierarchy used below imports the standard polarized
Thomson scattering matrix acting on Stokes brightness.  The scalar
$1+\cos^2\theta$ kernel displayed in Theorem~\ref{t:collision} is its
intensity projection, not a derivation of the $Q/U$ collision terms.  A
source-level OPH derivation of the polarized commit kernel from the recovered
vertex is not supplied here. The polarization sector and the associated $16/15$
diffusion factor are textbook imports rather than consequences of the displayed
operator.
\end{remark}

\begin{remark}[Blackbody provenance]
The Planck background spectrum is an initial-condition and thermal-history
import in this note, not a consequence of the displayed Thomson functional.
Even recoil-inclusive Comptonization conserves photon number and generically
approaches a Bose--Einstein spectrum with a chemical potential. Selecting
the zero-chemical-potential Planck spectrum requires the appropriate early
thermalization history, including photon-number-changing channels such as
double Compton scattering and bremsstrahlung, or an independently justified
blackbody initial condition. An OPH derivation of the blackbody spectrum
requires those channels and their rates rather than inferring
spectral uniqueness from angular detailed balance.
\end{remark}

\begin{remark}[Provenance of the normalization]
$\alpha_{\rm root}$ is an exact declared-map root. Its use as the physical
vertex strength is conditional on the three bridge premises in
Assumption~\ref{a:qed}. $m_e$ is an import
(Assumption~\ref{a:qed}): OPH supplies no source-derived charged-lepton
absolute mass, so $\sigma_T$'s normalization is
an explicit import. The recombination history $n_e(\tau)$ is imported under
Assumption~\ref{a:rec}.
\end{remark}

\begin{remark}[OPH-specific channels]\label{r:channels}
The declared two-channel continuation permits the following modifications
at linear order:
(i) a number-changing repair-exchange channel $\Gamma_{\rm rec}(k,a)$ adds a
monopole source and a spectral-distortion term; it enters $C[f]$
additively and requires an active fiber, a physical clock, and a
common-parent response;
(ii) a collar-packet stress $\rho_A,B_A(k,a)$ enters only the gravitational
source side (the Einstein constraints in \S\ref{s:hierarchy}), never $C[f]$.
Setting both declared corrections to zero recovers standard transport.
\end{remark}

\section{The Linear Multipole Hierarchy}\label{s:hierarchy}

\begin{theorem}[Reduction to the standard hierarchy]\label{t:hierarchy}
Linearize Theorems~\ref{t:liouville}--\ref{t:collision} on the background of
Assumption~\ref{a:bg} with scalar perturbations, expand the fractional
intensity brightness $F_\gamma(\mathbf k,\hat n,\tau)$ and polarization
brightness $G_\gamma$ as
$F_\gamma=\sum_\ell(-i)^\ell(2\ell+1)F_{\gamma\ell}P_\ell(\mu)$,
and likewise for $G_\gamma$, with $\mu=\hat k\cdot\hat n$. Writing
$\delta_\gamma=F_{\gamma0}$, $\theta_\gamma=\tfrac34 kF_{\gamma1}$,
$\sigma_\gamma=\tfrac12F_{\gamma2}$, and $\kdot\equiv a n_e\sigma_T$, the
system is (Ma--Bertschinger 1995, eq.~64):
\begin{align*}
\dot\delta_\gamma &= -\tfrac43\theta_\gamma+4\dot\phi,\\
\dot\theta_\gamma &= k^2\!\left(\tfrac14\delta_\gamma-\sigma_\gamma\right)
+k^2\psi+\kdot\,(\theta_b-\theta_\gamma),\\
2\dot\sigma_\gamma &= \tfrac{8}{15}\theta_\gamma-\tfrac35 kF_{\gamma3}
-\tfrac95\kdot\,\sigma_\gamma
+\tfrac{1}{10}\kdot\,(G_{\gamma0}+G_{\gamma2}),\\
\dot F_{\gamma\ell} &= \frac{k}{2\ell+1}\left[\ell F_{\gamma(\ell-1)}
-(\ell+1)F_{\gamma(\ell+1)}\right]-\kdot F_{\gamma\ell},
\qquad \ell\ge3,\\
\dot G_{\gamma\ell} &= \frac{k}{2\ell+1}\left[\ell G_{\gamma(\ell-1)}
-(\ell+1)G_{\gamma(\ell+1)}\right]
+\kdot\left[-G_{\gamma\ell}
+\tfrac12\left(F_{\gamma2}+G_{\gamma0}+G_{\gamma2}\right)
\left(\delta_{\ell0}+\tfrac{\delta_{\ell2}}{5}\right)\right],
\end{align*}
coupled to baryons,
\[
\dot\delta_b=-\theta_b+3\dot\phi,\qquad
\dot\theta_b=-\ah\theta_b+c_b^2k^2\delta_b+k^2\psi
+\frac{4\bar\rho_\gamma}{3\bar\rho_b}\,\kdot\,(\theta_\gamma-\theta_b),
\]
and to the Einstein constraints of Assumption~\ref{a:geom},
\[
k^2\phi+3\ah\!\left(\dot\phi+\ah\psi\right)=-4\pi G a^2\,\delta\rho_{\rm tot},
\qquad
k^2(\phi-\psi)=12\pi G a^2(\bar\rho+\bar p)\sigma_{\rm tot},
\]
where $\delta\rho_{\rm tot}$ and $\sigma_{\rm tot}$ include the gated
$\rho_A,B_A$ contributions of Remark~\ref{r:channels}. These contributions are
zero on the receipt-failing branch. With $\Gamma_{\rm rec}=B_A=0$ this is verbatim the standard
Einstein--Boltzmann system.
\end{theorem}

\begin{proof}
For scalar intensity, insert the geodesic characteristics of
Theorem~\ref{t:liouville} on the perturbed metric, the collision functional
of Theorem~\ref{t:collision}, keep first order, and project on Legendre
multipoles. The $G_{\gamma\ell}$ equations additionally use the imported
polarized Thomson matrix stated above. The subsequent calculation is the
textbook computation (Ma--Bertschinger 1995; Dodelson, \emph{Modern
Cosmology}). The two conditional OPH channels appear at their unique insertion
points.
\end{proof}

\section{Acoustic normal form, damping, and the line of sight}

\begin{proposition}[Acoustic oscillations]\label{p:acoustic}
In the tight-coupling regime $\kdot\gg k,\ \ah$, neglect baryon thermal
pressure and define
$R\equiv 3\bar\rho_b/4\bar\rho_\gamma$ and $\Theta_0=\delta_\gamma/4$. To
leading order in $1/\kdot$,
\[
\frac{\Dd}{\Dd\tau}\!\left[(1+R)\,\dot\Theta_0\right]
+\frac{k^2}{3}\,\Theta_0
=-\frac{k^2}{3}(1+R)\,\psi
+\frac{\Dd}{\Dd\tau}\!\left[(1+R)\,\dot\phi\right],
\]
a driven oscillator with sound speed $c_s^2=1/[3(1+R)]$. Its homogeneous
WKB phase, when the coefficients vary slowly compared with an oscillation,
is $k r_s$, where
\[
r_s(\tau_\ast)=\int_0^{\tau_\ast}\!c_s\,\Dd\tau,
\]
at last scattering $\tau_\ast$. A cosine initial phase gives the approximate
unforced extrema $k_n r_s\simeq n\pi$. Baryon loading changes the
equilibrium offset and relative compression and rarefaction amplitudes;
evolving gravitational potentials drive the oscillator. Initial phases,
free-streaming species, finite visibility and angular projection shift the
observed peak positions and heights. The oscillator therefore identifies the
acoustic mechanism without fixing exact peaks of a measured TT spectrum.
\end{proposition}

\begin{proposition}[Silk damping]\label{p:damping}
At next order in $1/\kdot$ the oscillator acquires the diffusion envelope
$e^{-k^2/k_D^2}$ with
\[
k_D^{-2}=\int_0^{\tau}\frac{\Dd\tau'}{6(1+R)\,\kdot}
\left[\frac{R^2}{1+R}+\frac{16}{15}\right],
\]
the $16/15$ arising from the polarization terms of
Theorem~\ref{t:hierarchy}; photon diffusion between commits is the physical
mechanism, and its OPH reading is transparent: on scales crossed by a random
walk of commit-free transport legs, the record field cannot maintain
contrast.
\end{proposition}

\begin{proposition}[Line-of-sight solution]\label{p:los}
For $k>0$, with visibility $g(\tau)=\kdot\,e^{-\kappa(\tau,\tau_0)}$ and an
optically thick initial endpoint, the temperature multipoles today, for
$\ell\geq2$, are
\[
\Theta_\ell(k,\tau_0)=\int_0^{\tau_0}\!\Dd\tau\;
\Big\{
g\left[\Theta_0+\psi+\tfrac{\Pi}{4}\right]
j_\ell\big(k(\tau_0-\tau)\big)
+\Big(\frac{g\,v_b}{k}\Big)^{\!\cdot} j_\ell
+\frac{3}{4k^2}\big(g\Pi\big)^{\!\cdot\cdot}j_\ell
+e^{-\kappa}\big(\dot\psi+\dot\phi\big)\,j_\ell
\Big\},
\]
($\Theta_\ell=F_{\gamma\ell}/4$,
$\Pi=(F_{\gamma2}+G_{\gamma0}+G_{\gamma2})/4$,
$v_b=\theta_b/k$; every unlabelled
$j_\ell$ has argument $k(\tau_0-\tau)$), and
\[
C_\ell^{TT}=4\pi\!\int\!\frac{\Dd k}{k}\,
\Delta_\zeta^2(k)\,\left|\frac{\Theta_\ell(k,\tau_0)}{\zeta(k)}\right|^2 .
\]
Given the conditional $(A_s,n_s)$ of Assumption~\ref{a:init} and the imports of
Assumptions~\ref{a:qed}--\ref{a:rec}, this is the standard linear transfer
formula. Agreement of an independent implementation with CAMB at identical
inputs and sub-percent tolerances is a separate numerical acceptance test.
\end{proposition}

The normalization can be checked directly from the stated hierarchy.
Dividing the intensity collision term of Ma--Bertschinger by four gives
$-\kdot\Pi P_2(\mu)/2$ and the Doppler term $-i\mu\kdot\theta_b/k$.
Along the line of sight, $P_2(\mu)=(3\mu^2-1)/2$ and
$\partial_\tau e^{ik\mu(\tau-\tau_0)}=ik\mu e^{ik\mu(\tau-\tau_0)}$.
Integration by parts therefore gives $g\Pi/4+3(g\Pi)^{\cdot\cdot}/(4k^2)$
and $(g\theta_b)^{\cdot}/k^2$, respectively. The initial boundary terms
vanish by optical thickness, and the observer terms affect only the
monopole and dipole. These conventions specify the signs and factors in
the Seljak--Zaldarriaga line-of-sight construction.

\Needspace{12\baselineskip}
\section{Premises and Numerical Tests}\label{s:premises}

\begin{longtable}{@{}L{0.16\linewidth}L{0.5\linewidth}L{0.26\linewidth}@{}}
\toprule
\textbf{Condition} & \textbf{Statement} & \textbf{Classification}\\
\midrule
Photon refinement convergence &
Finite-lattice refinement Liouville: coarse-grained photon-class record
transport has quotient lumpability, tightness, convergent initial laws and
first conditional moments, vanishing scaled second conditional moments,
generator convergence on a frozen test basis, unique geodesic characteristics,
and vanishing schedule, refinement, boundary, mass-shell, and measure defects.
This condition is independent
of repair confluence and of the orientation-doubled label count. &
scalar free-flow generator supplied on the flat declared family
(Remark~\ref{r:scalar-continuum}); vector, record-process and
curved-background clauses: simulator test below\\
Thomson normalization &
$\sigma_T$ normalization: absolute $m_e$ is a declared import; the
charged-family absolute-scale theorem is not supplied. &
declared import\\
Recombination history &
$x_e(\tau)$ recombination history from imported atomic data; contains no
CMB-target information. &
declared import\\
Numerical transfer test &
Independent hierarchy implementation matches CAMB temperature and polarization spectra at stated
tolerances with identical inputs. &
not supplied; engineering test specified below\\
Tilt certificate &
Conditional \(n_s=1-P_\star/48\); source generator, weighted half, and clock
normalization are not supplied. &
not supplied\\
Amplitude + lift &
$A_s$ and radial screen-to-$k$ map. &
conditional radial theorem; finite source receipt not supplied\\
Kernels &
$\Gamma_{\rm rec}$, $\rho_A$, $B_A$ physical receipts. &
not supplied\\
Likelihood &
Official likelihood, provenance, and data-use contract. &
not supplied\\
\bottomrule
\end{longtable}

\textbf{Simulator test for photon refinement convergence.} Freeze a smooth compactly
supported test basis. At each regulator estimate the first conditional moment
$\widehat L_r\phi$ and scaled conditional variance
$\widehat\Gamma_r(\phi)$, together with quotient-measure, mass-shell,
schedule, refinement, boundary-leakage, and collision/free-channel residuals.
Pass requires $\widehat L_r\phi\to X_H\phi$ and
$\widehat\Gamma_r(\phi)\to0$ on the frozen basis with the preregistered
uniform compact envelopes. Random-schedule tests are supplemental.
A nonzero second-moment limit excludes the deterministic Liouville theorem;
the limiting jump sizes and generator distinguish diffusion from jump transport.

\paragraph{Native repair diffusion and the photon interface.}
For the specified fair conservative port repair, mean load evolves as
\(e^{-tL/2}\), with \(t\) the native attempt count divided by the seam
count, and stationary density has covariance
\(\kappa e^{-|t-s|L/2}\Pi\). Its integrated readback has a
fixed-graph Gaussian limit with inverse-Laplacian covariance; the literal
finite-time average has the different normalization
\(\kappa h(\lambda T/2)\). The exact derivation, local attempted-event
record correction and finite-state controls are given in the
screen-microphysics companion~\cite{ophmicro} and preserved in the
observer-dynamics evidence archive~\cite{ophobsdyn}.
Along one common subsequence of refining source screens, actual native
density snapshots have joint Gaussian limits with covariance
\(\exp(-|\tau-u|A)\), where \(\tau\) is the native attempt count divided
by \(1260\cdot16^{l}\) and \(A\) is a positive self-adjoint operator of
heat-trace spectral dimension two on the
screen~\cite{ophtimefields,ophcontinuum}. Coherent waves
\(\cos(t\sqrt A)\) require a declared wave processor whose parameter differs
from passive waiting, and they propagate at most at speed \(1/\sqrt2\) in the
flat comparison distance per coherent wave parameter~\cite{ophmodel}. This two-dimensional screen
operator supplies no \(3+1\) null-geodesic flow or photon sector, and its
density limit is a stationary Gaussian process, not a deterministic
Liouville flow.
These statements instantiate a diffusion/record mechanism. They do not
verify the vanishing conditional second moment required for a deterministic
photon Liouville limit, nor do they provide polarization, Thomson
normalization or recombination.

A nondecreasing global change of time keeps every positive mean mode
\(\exp[-\lambda g(t)/2]\) positive and decreasing. It cannot supply an
oscillatory wave propagator. For any current-state centered observable of
a stationary reversible continuous-time repair chain, its autocorrelation
is a nonnegative sum of decaying exponentials. The corresponding
discrete-time conclusion requires a nonnegative spectrum, as in a lazy
chain; a nonlazy two-state flip is a counterexample. These restrictions
apply to the stated generator and readout, not to extended dynamics with
momentum, active forcing, nonstationary preparation or different records.
The passive source matching and acoustic initial data used by a conditional
CMB transport calculation are additional physical inputs.

\section{Scope of the Conditional Reduction}
The cited branch theorems supply the \emph{scalar angular form} of the
transport and locate the two conditional OPH kernels. The polarization
hierarchy, blackbody background, absolute normalizations, and recombination
history are imported. A physical cosmic microwave background prediction
requires a source-level polarized collision bridge, blackbody and
thermalization provenance, photon refinement convergence, the numerical
transfer test, a physical radial lift, a finite primordial source,
active-kernel evidence, and an immutable likelihood contract. This reduction
certifies no initial conditions, normalizations, recombination history, or
likelihood, and it reads no Planck data. Under the anti-fitting rule, the
acoustic peak structure is attributable to standard plasma physics. OPH's
falsifiable exposure is concentrated in the primordial inputs.

\begin{thebibliography}{99}

\bibitem{mabertschinger}
C.-P. Ma and E. Bertschinger, \emph{Cosmological Perturbation Theory
in the Synchronous and Conformal Newtonian Gauges}, ApJ 455, 7 (1995).

\bibitem{dodelson}
S. Dodelson and F. Schmidt, \emph{Modern Cosmology}, 2nd ed., Academic
Press (2020).

\bibitem{seljakzaldarriaga}
U. Seljak and M. Zaldarriaga, \emph{A Line-of-Sight Integration
Approach to Cosmic Microwave Background Anisotropies}, ApJ 469, 437 (1996).

\bibitem{lionspaul}
P.-L. Lions and T. Paul, \emph{Sur les mesures de Wigner}, Rev. Mat.
Iberoam. 9, 553 (1993).

\bibitem{bashinskyseljak}
S. Bashinsky and U. Seljak, \emph{Signatures of Relativistic Neutrinos
in CMB Anisotropy and Matter Clustering}, Phys. Rev. D 69, 083002 (2004).

\bibitem{huwhite}
W. Hu and M. White, \emph{CMB Anisotropies: Total Angular Momentum
Method}, Phys. Rev. D 56, 596 (1997).

\bibitem{ophobservers}
B.~M\"uller et al., \emph{Observers Are All You Need}.
\url{https://github.com/FloatingPragma/observer-patch-holography/blob/r2056/paper/observers_are_all_you_need.pdf}.

\bibitem{ophconsensus}
B.~M\"uller et al., \emph{Reality as a Consensus Protocol}.
\url{https://github.com/FloatingPragma/observer-patch-holography/blob/r2056/paper/reality_as_consensus_protocol.pdf}.

\bibitem{ophEinstein}
B.~M\"uller et al., \emph{Recovering Observer Spacetime and Einstein
Dynamics from Overlap Consistency}.
\url{https://github.com/FloatingPragma/observer-patch-holography/blob/r2056/paper/recovering_observer_spacetime_and_einstein_dynamics_from_overlap_consistency.pdf}.

\bibitem{ophmicro}
B.~M\"uller et al., \emph{Twelve-Port Icosahedral Carriers and Native Seam Repair: The Carrier Layer of Observer-Patch Holography}.
\url{https://github.com/FloatingPragma/observer-patch-holography/blob/r2056/paper/screen_microphysics_and_observer_synchronization.pdf}.

\bibitem{ophcmb}
B.~M\"uller, \emph{Finite-Source Cosmology from Observer-Patch Holography:
Source, Transfer, and Likelihood Boundaries}.
\url{https://github.com/FloatingPragma/observer-patch-holography/blob/r2056/cosmology/oph_cosmology_finite_source_cmb_program.pdf}.

\bibitem{ophflagship}
B.~M\"uller et al., \emph{Observers Who Read Their Universe into Existence:
Finite Consensus, an Executable Observer, and Standard Physics}.
\url{https://github.com/FloatingPragma/observer-patch-holography/blob/r2056/flagship/from_observer_consensus_to_standard_physics.pdf}.

\bibitem{ophspine}
B.~Mueller, M.~Poneder, J.~Kim and D.~A.~Friedman (2026). \emph{An Executable Observer for
Observer Patch Holography: Specification, Executed Components and Ordered Joins}.
\url{https://github.com/FloatingPragma/observer-patch-holography/blob/r2056/paper/executable_observer_construction.pdf}.

\bibitem{ophmodel}
FloatingPragma (2026). The runnable OPH model: observer implementation,
tests, receipts and construction guide.
\url{https://github.com/FloatingPragma/observer-patch-holography/tree/r2056/model}.

\bibitem{ophcontinuum}
FloatingPragma (2026). A spatial fluctuation field from one native history:
the common spatial generator, its locality and spectral dimension, and the
boundary-trace test.
\url{https://github.com/FloatingPragma/observer-patch-holography/blob/r2056/model/physical_space_time/OCCUPATION_FIELD.md}.

\bibitem{ophtimefields}
FloatingPragma (2026). Joint time and field laws on matched source screens.
\url{https://github.com/FloatingPragma/observer-patch-holography/blob/r2056/model/observer_time_fields/REFINEMENT.md}.

\bibitem{ophentropy}
FloatingPragma (2026). Tangent first laws and finite entropy changes:
native reference entropy and region boundaries.
\url{https://github.com/FloatingPragma/observer-patch-holography/blob/r2056/code/geometry/ENTROPY_FIRST_LAW_AUDIT.md}.

\bibitem{ophcausality}
FloatingPragma (2026). Exact native causal cones.
\url{https://github.com/FloatingPragma/observer-patch-holography/blob/r2056/model/common_source/CAUSALITY.md}.

\bibitem{ophpublicdata}
FloatingPragma (2026). OPH predictions and public-data comparisons.
\url{https://github.com/FloatingPragma/observer-patch-holography/blob/r2056/docs/research/PUBLIC_DATA_PREDICTION_AUDIT.md}.

\bibitem{ophobsdyn}
FloatingPragma (2026). Observer-dynamics analysis archive: derivation
receipts and finite-state controls.
\url{https://github.com/FloatingPragma/observer-patch-holography/tree/r2056/evidence/observer_dynamics_20260925/sim-analysis/codex}.

\end{thebibliography}

\section*{AI Assistance Disclosure}
This research project used research-grade commercial models, including
Anthropic's Fable and OpenAI's GPT-6 Astra, for research support, software
development, editing, and synthesis. The authors are responsible for the
paper's claims, methods, and final text.

\end{document}
